\documentclass[11pt]{article}
\usepackage{amsmath,amssymb,amsthm,geometry,booktabs,hyperref}
\geometry{margin=1in}
\newtheorem{theorem}{Theorem}
\newtheorem{lemma}{Lemma}
\newcommand{\Tr}{\operatorname{Tr}}
\begin{document}
\title{An auxiliary-spin bound for second-order anticoherent subspaces}
\author{Independent catalogue calculation}
\date{30 October 2026}
\maketitle
\begin{theorem}
For every integer $N\ge8$, a complex second-order anticoherent subspace
of $\operatorname{Sym}^N(\mathbb C^2)$ has dimension
\[
 \boxed{\quad k\le\left\lfloor\frac{N-2}{4}\right\rfloor.\quad}
\]
In particular, $k_{\max}(13,2)=2$.
\end{theorem}
The upper bound applies to arbitrary complex subspaces. It assumes neither
Dicke sparsity nor a rotational symmetry of the subspace or its states.
The auxiliary spin below is a proof device, not an additional physical
qubit included in the definition of the code.

\section*{1. Whole-subspace second moments}
Let $j=N/2$, $\mathcal H_j\cong\mathbb C^{N+1}$, and let $P$ project
onto a $k$-dimensional second-order AC subspace $\mathcal C$. Polarization
of the AC conditions gives
\[
 PJ_aP=0,\qquad PJ_aJ_bP=c\delta_{ab}P,\qquad
 c=\frac{j(j+1)}3=\frac{N(N+2)}{12}.
\]
The symmetric second moments give the symmetric part of the second
identity. Its antisymmetric part vanishes because
$[J_a,J_b]=i\epsilon_{abc}J_c$ and the first identity holds. Thus all
ordered products, including the cross terms between different code
vectors, are covered.

Introduce an auxiliary spin $s=3/2$, with operators $S_a$ on
$\mathbb C^4$. Set
\[
 X=\sum_{a=x,y,z}J_a\otimes S_a,
 \qquad \Pi=P\otimes I_4.
\]
Then
\[
 \Pi X\Pi=0,\qquad
 \Pi X^2\Pi=cP\otimes\sum_aS_a^2
 =\frac{5N(N+2)}{16}\Pi.
\]
Consequently the Hermitian operator
\[
 \boxed{W=X^2-\frac{N-4}{2}X-\frac{5N(N+2)}{16}I}
\]
satisfies
\[
 \Pi W\Pi=0.
\]
This means that $\mathcal C\otimes\mathbb C^4$, a complex subspace of
dimension $4k$, is totally isotropic for the Hermitian form defined by $W$.

\section*{2. The spectrum has only $N-2$ positive directions}
The tensor product of spins $j$ and $3/2$ decomposes into total spins
$J=j-3/2,j-1/2,j+1/2,j+3/2$. Since
\[
 X=\frac12\bigl(J_{\mathrm{tot}}^2-j(j+1)I-s(s+1)I\bigr),
\]
its four eigenvalues and the corresponding eigenvalues of $W$ are
\[
\begin{array}{c|c|c|c}
J&\text{multiplicity}&X&W\\\hline
j-3/2&N-2&-3(N+2)/4&(N+2)(5N-3)/8\\
j-1/2&N&-(N+8)/4&-N(N-7)/8\\
j+1/2&N+2&(N-6)/4&-(3N^2+N+6)/8\\
j+3/2&N+4&3N/4&-N(N-7)/8.
\end{array}
\]
For $N\ge8$, the first is strictly positive and the remaining three are
strictly negative. There are no zero eigenvalues. Therefore the inertia
of $W$ is
\[
 (n_+,n_-,n_0)=(N-2,3N+6,0).
\]

\begin{lemma}
If a nonsingular Hermitian form has $p$ positive and $q$ negative
directions, every totally isotropic complex subspace has dimension at
most $\min(p,q)$.
\end{lemma}
\begin{proof}
Projection of an isotropic subspace to the positive eigenspace is
injective: a vector in its kernel would lie entirely in the negative
eigenspace and could have zero quadratic form only if it were zero.
The same argument applies to the negative eigenspace.
\end{proof}
Applying the lemma to $\mathcal C\otimes\mathbb C^4$ gives
\[
 4k\le N-2,
\]
which proves the general theorem. At $N=7$ two eigenvalues of $W$ vanish,
so this nonsingular argument does not apply. This is consistent with the
known two-dimensional second-order AC subspace at seven qubits.

\section*{3. The thirteen-qubit case and exact attainment}
At $N=13$,
\[
 W=X^2-\frac92X-\frac{975}{16}I.
\]
Its eigenvalues are $465/4,-39/4,-263/4,-39/4$, on total spins
$5,6,7,8$ of dimensions $11,13,15,17$, respectively. Hence $4k\le11$
and $k\le2$.

In normalized Dicke notation $D_r=|D_{13}^{(r)}\rangle$, an attaining
orthonormal basis is
\[
 u=\frac{\sqrt{10}D_0+\sqrt{26}D_9}{6},\qquad
 v=\frac{-\sqrt{26}D_4+\sqrt{10}D_{13}}{6}.
\]
Every support gap exceeds two. All rank-one/rank-two ladder terms
therefore vanish, including cross terms between the two vectors. Each
vector has zero magnetic mean and second moment $65/4$, establishing
the complete second-order AC compression identities. Thus
\[
 \boxed{k_{\max}(13,2)=2.}
\]
The earlier one-parameter phase family also attains this maximum; the
upper bound is independent of that family and does not assume that a
hypothetical code contains either displayed vector.

\section*{4. Reproducible independent verification}
The accompanying \texttt{verify.py} uses Python's standard library only.
It reconstructs the physical spin operators in rescaled Dicke bases
$b_r=\sqrt{\binom Nr}D_r$ and
$c_a=\sqrt{\binom3a}|D_3^{(a)}\rangle$:
\[
 J_-b_r=(r+1)b_{r+1},\qquad J_+b_r=(N-r+1)b_{r-1},
\]
with analogous auxiliary-spin actions. The resulting matrices have
rational entries and are self-adjoint for the explicitly retained
positive diagonal Gram matrix.

The verifier builds every conserved-weight block of the physical $X$.
It constructs rational spectral projectors as polynomials in those
blocks and verifies their idempotence, orthogonality, completeness,
weighted self-adjointness, exact eigenvalues and multiplicities. Checks
at $N=7,8,9,12,13,14,16,17$ include the exceptional seven-qubit nullspace
and the signature claimed for $N\ge8$.

For two members of the thirteen-qubit family, the verifier independently
checks all 32 rank-one/rank-two cross-compression entries and all 64
entries of the direct physical auxiliary-spin compression
$(U\otimes I_4)^\dagger W(U\otimes I_4)=0$, using exact arithmetic in
$\mathbb Q(\sqrt{11})$. No numerical optimizer, Wigner-symbol package,
or symbolic algebra package is required. These finite checks validate
the implementation and the explicit witnesses; the formula and inertia
argument above establish the theorem for all $N\ge8$.
\end{document}
